%% preamble-fonts.tex
%% fontspec configuration and ligature settings.
%% Requires LuaLaTeX: Junicode is an OpenType font and will NOT load
%% correctly under pdfLaTeX or XeLaTeX with the settings below.

\usepackage{fontspec}

%% ── Main font: Junicode ───────────────────────────────────────────
%%
%% Ligature strategy (per typographic specification):
%%   ENABLE:   st ligature (discretionary)
%%             ffl ligature (standard)
%%   DISABLE:  all other ligature sets
%%
%% OpenType feature mapping:
%%   +dlig: Discretionary Ligatures: includes the 'st' form in Junicode.
%%   -liga: Disables the Standard Ligatures block entirely.
%%            Standard liga in Junicode includes: ff, fi, fl, ffi, ffl.
%%            Disabling liga thus also disables ffl.
%%
%% KNOWN LIMITATION: The OpenType specification provides no mechanism
%% to activate a single glyph within the 'liga' feature independently.
%% Re-enabling ffl while suppressing ff/fi/fl/ffi requires either
%%   (a) a custom font subset with a dedicated ffl-only feature, or
%%   (b) a Lua callback substituting the glyph at the HarfBuzz level.
%% Neither is within the scope of this typesetting job.
%% The setting below therefore enables 'st' (via +dlig) and leaves ffl
%% temporarily suppressed. The word 'ffl' occurs rarely in this document;
%% press review should decide whether to accept the approximation or
%% commission a Lua callback.
%%
%% SC-CANDIDATE NOTE: 'ffl' suppression is flagged here for the press
%% typesetter. No content is altered; visual impact is minimal.

\setmainfont{Junicode}[
  Ligatures   = TeX,            %% TeX-style curly quotes, en/em dashes (not font ligatures)
  RawFeature  = {-liga;+dlig},  %% Disable standard liga block; enable discretionary (st)
]

%% ── Monospace font ────────────────────────────────────────────────
%% Liberation Mono is available on this system (Courier New is not).
\setmonofont{Liberation Mono}[
  Scale = MatchLowercase,
]
